\sectionheading{5.29 Python-এ Code লেখার মাধ্যম}

Python Program বিভিন্ন ধরনের Environment বা Tool ব্যবহার করে
লেখা ও Execute করা যায়।

কাজের ধরন অনুযায়ী কখনো সরাসরি Instruction লিখে সঙ্গে সঙ্গে
Result দেখা যায়, আবার কখনো Program File তৈরি করে Code
Save ও Run করা হয়।

%================================================
% Main Ways
%================================================

\booksubheading{Python Code লেখার প্রধান মাধ্যম}

Python Code লেখার পদ্ধতি সহজভাবে প্রধানত দুইভাবে বোঝা যায়—

\begin{enumerate}

\item \textbf{Interactive Mode / Python Shell}

\item \textbf{Script Mode / Editor}

\end{enumerate}

এ ছাড়া Code লেখাকে সহজ করার জন্য বিভিন্ন
\textbf{IDE, Code Editor, Notebook ও Online Environment}
ব্যবহার করা যায়।

\begin{center}

\textbf{Interactive Mode}
\qquad
\textbf{ও}
\qquad
\textbf{Script Mode}

\end{center}

%================================================
% Interactive Mode
%================================================

\replaceableimage[60mm]
{images/python-shell.png}
{Python IDLE Shell-এ Prompt, Command ও Output}

\begin{center}
\footnotesize\textbf{চিত্র: Python Interactive Shell-এ সরাসরি Code Execute করা}
\end{center}

\booksubheading{১. Interactive Mode / Python Shell}

Interactive Mode-এ Python Shell-এর Prompt-এ সরাসরি
Instruction লিখে Enter চাপলে সাধারণত সঙ্গে সঙ্গে Result দেখা যায়।

Python Shell-এর Prompt হলো—

\begin{codebox}

\begin{english}
\ttfamily
>>>
\end{english}

\end{codebox}

\begin{examplebox}{উদাহরণ}

\begin{codebox}

\begin{english}
\ttfamily
>>> 10 + 5\\
15
\end{english}

\end{codebox}

এখানে
\textenglish{\ttfamily 10 + 5}
লিখে Enter চাপার সঙ্গে সঙ্গে
\textenglish{\ttfamily 15}
Result পাওয়া গেছে।

\end{examplebox}

\textbf{Interactive Mode-এর বৈশিষ্ট্য—}

\begin{itemize}

\item একবারে ছোট Instruction পরীক্ষা করা যায়।

\item Result সঙ্গে সঙ্গে দেখা যায়।

\item ছোট Calculation ও Code Test-এর জন্য সুবিধাজনক।

\item প্রতিবার Program File তৈরি করা প্রয়োজন হয় না।

\item বড় Program লেখার জন্য তুলনামূলক কম সুবিধাজনক।

\end{itemize}

\begin{rememberbox}

\textbf{\color{deepOrange}সহজভাবে মনে রাখবে}

\begin{center}

\textbf{Command লিখি}
\(\longrightarrow\)
\textbf{Enter চাপি}
\(\longrightarrow\)
\textbf{সঙ্গে সঙ্গে Result}

\end{center}

\end{rememberbox}

%================================================
% Script Mode
%================================================

\booksubheading{২. Script Mode / Editor}

Script Mode-এ একাধিক Python Statement একটি Program File-এ
লেখা হয়।

File-টি
\textbf{\textenglish{\ttfamily .py}}
Extension-এ Save করে পরে পুরো Program একসঙ্গে Execute করা যায়।

\begin{examplebox}{উদাহরণ}

\begin{codebox}

\begin{english}
\ttfamily
a = 10\\
b = 20\\
sum = a + b\\
print(sum)
\end{english}

\end{codebox}

Program File-এর নাম হতে পারে—

\begin{codebox}

\begin{english}
\ttfamily
sum.py
\end{english}

\end{codebox}

\end{examplebox}

\textbf{Script Mode-এর বৈশিষ্ট্য—}

\begin{itemize}

\item একাধিক Statement নিয়ে Program লেখা যায়।

\item Code Program File হিসেবে Save করা যায়।

\item প্রয়োজন হলে পরে Code Edit করা যায়।

\item একই Program বারবার Run করা যায়।

\item বড় ও সংরক্ষণযোগ্য Program তৈরির জন্য উপযোগী।

\end{itemize}

\begin{rememberbox}

\textbf{\color{deepOrange}সহজভাবে মনে রাখবে}

\begin{center}

\textbf{Code লিখি}
\(\longrightarrow\)
\textbf{\textenglish{\ttfamily .py} File Save করি}
\(\longrightarrow\)
\textbf{Run করি}
\(\longrightarrow\)
\textbf{Output পাই}

\end{center}

\end{rememberbox}

%================================================
% Interactive vs Script
%================================================

\booksubheading{Interactive Mode ও Script Mode-এর পার্থক্য}

\begin{tabularx}
{\linewidth}
{
|>{\bfseries\RaggedRight\arraybackslash}p{0.20\linewidth}
|Y
|Y|
}

\hline

\rowcolor{softBlue}

\multicolumn{1}{|c|}{\textbf{বিষয়}}
&
\multicolumn{1}{c|}{\textbf{Interactive Mode}}
&
\multicolumn{1}{c|}{\textbf{Script Mode}}

\\

\hline

Code লেখা
&
Shell Prompt-এ সরাসরি লেখা হয়।
&
Editor বা Program File-এ লেখা হয়।
\\

\hline

Execution
&
Instruction দেওয়ার সঙ্গে সঙ্গে Execute হয়।
&
Program File Run করলে Execute হয়।
\\

\hline

Result
&
সাধারণত সঙ্গে সঙ্গে দেখা যায়।
&
Program Run করার পর দেখা যায়।
\\

\hline

File Save
&
সাধারণত বাধ্যতামূলক নয়।
&
\textenglish{\ttfamily .py}
File হিসেবে Save করা হয়।
\\

\hline

ব্যবহার
&
ছোট Test ও Calculation-এর জন্য সুবিধাজনক।
&
পূর্ণ Program লেখার জন্য সুবিধাজনক।
\\

\hline

পুনঃব্যবহার
&
আগের Instruction সংরক্ষণ ও পুনঃব্যবহার তুলনামূলক কম সুবিধাজনক।
&
Save করা Program সহজে আবার Run ও Edit করা যায়।
\\

\hline

\end{tabularx}

\begin{rememberbox}

\textbf{\color{deepOrange}পরীক্ষায় গুরুত্বপূর্ণ}

\textbf{Interactive Mode}
\(\longrightarrow\)
Immediate Result

\medskip

\textbf{Script Mode}
\(\longrightarrow\)
Saved Program

\end{rememberbox}

%================================================
% IDLE
%================================================

\booksubheading{IDLE-এ Python Code লেখা}

Python-এর পরিচিত Environment
\textbf{IDLE}-এ সাধারণভাবে দুইভাবে কাজ করা যায়—

\textbf{ক. Python Shell}

Shell-এ সরাসরি Instruction লিখে Result দেখা যায়।

\begin{codebox}

\begin{english}
\ttfamily
>>> print("Hello")\\
Hello
\end{english}

\end{codebox}

\textbf{খ. Editor Window}

Editor Window-তে একাধিক Statement লিখে Program File হিসেবে
Save ও Run করা যায়।

\textbf{সাধারণ ধাপ—}

\begin{center}

\textbf{New File}
\(\longrightarrow\)
\textbf{Code লেখা}
\(\longrightarrow\)
\textbf{Save}
\(\longrightarrow\)
\textbf{Run Module}
\(\longrightarrow\)
\textbf{Output}

\end{center}

Python Program File সাধারণত
\textbf{\textenglish{\ttfamily .py}}
Extension-এ Save করা হয়।

%================================================
% Code Editor
%================================================

\booksubheading{Code Editor ব্যবহার করে Python লেখা}

Python Program লেখার জন্য বিভিন্ন Code Editor ব্যবহার করা যায়।

উদাহরণ—

\begin{center}

\textbf{Visual Studio Code (VS Code)}

\end{center}

Code Editor-এ সাধারণত—

\begin{compacttwocolpoints}

\item Program Code লেখা যায়।

\item Code Edit করা যায়।

\item Syntax Highlighting পাওয়া যায়।

\item Program File Save করা যায়।

\item প্রয়োজনীয় Setup থাকলে Program Run করা যায়।

\end{compacttwocolpoints}

\begin{rememberbox}

\textbf{\color{deepOrange}মনে রাখবে}

Code Editor-এর মূল কাজ—

\begin{center}

\textbf{Write}
\(\longrightarrow\)
\textbf{Edit}
\(\longrightarrow\)
\textbf{Save Code}

\end{center}

\end{rememberbox}

%================================================
% IDE
%================================================

\booksubheading{IDE ব্যবহার করে Python লেখা}

IDE-তে Code লেখা, Run, Debug এবং Project Management-এর মতো
একাধিক Development সুবিধা একই Environment-এ পাওয়া যায়।

Python-এর পরিচিত IDE-এর উদাহরণ—

\begin{itemize}

\item \textbf{IDLE}

\item \textbf{PyCharm}

\end{itemize}

IDE-তে সাধারণত—

\begin{center}

\textbf{Code Editor}
+
\textbf{Run}
+
\textbf{Debug}
+
\textbf{Project Tool}

\end{center}

সুবিধা পাওয়া যায়।

%================================================
% Jupyter Notebook
%================================================

\booksubheading{Jupyter Notebook}

Python Code লেখার আরেকটি পরিচিত মাধ্যম হলো
\textbf{Jupyter Notebook}।

এখানে Code সাধারণত ছোট ছোট
\textbf{Cell}-এ লেখা ও Execute করা যায়।

একটি Code Cell Run করলে তার Output সাধারণত সেই Cell-এর
কাছাকাছি দেখা যায়।

\begin{center}

\textbf{Cell-এ Code}
\(\longrightarrow\)
\textbf{Run}
\(\longrightarrow\)
\textbf{Output}

\end{center}

Jupyter Notebook বিশেষভাবে—

\begin{itemize}

\item Data Analysis

\item Scientific Computing

\item Learning ও Experiment

\item Machine Learning

\end{itemize}

জাতীয় কাজে ব্যবহৃত হয়।

%================================================
% Online Environment
%================================================

\booksubheading{Online Python Environment}

কিছু Online Programming Environment ব্যবহার করে Browser থেকেই
Python Code লেখা ও Run করা যায়।

এক্ষেত্রে Local Computer-এ আলাদাভাবে Python Install না করেও
Practice করা সম্ভব হতে পারে।

\begin{center}

\textbf{Browser}
\(\longrightarrow\)
\textbf{Code লিখি}
\(\longrightarrow\)
\textbf{Run}
\(\longrightarrow\)
\textbf{Output}

\end{center}

\textbf{সুবিধা—}

\begin{itemize}

\item Installation ছাড়াই Practice করা যায়।

\item বিভিন্ন Device থেকে ব্যবহার করা যায়।

\item Beginner-এর Practice-এর জন্য সুবিধাজনক।

\end{itemize}

\textbf{সীমাবদ্ধতা—}

\begin{itemize}

\item সাধারণত Internet প্রয়োজন হয়।

\item Local Computer-এর সব সুবিধা পাওয়া নাও যেতে পারে।

\item Platform অনুযায়ী Feature ভিন্ন হতে পারে।

\end{itemize}

%================================================
% Which Mode to Use
%================================================

\booksubheading{কোন কাজের জন্য কোন মাধ্যম?}

\begin{tabularx}
{\linewidth}
{
|>{\RaggedRight\arraybackslash}p{0.42\linewidth}
|Y|
}

\hline

\rowcolor{softBlue}

\multicolumn{1}{|c|}{\textbf{প্রয়োজন}}
&
\multicolumn{1}{c|}{\textbf{উপযোগী মাধ্যম}}

\\

\hline

ছোট Calculation
&
Interactive Shell
\\

\hline

ছোট Code পরীক্ষা
&
Interactive Mode
\\

\hline

পূর্ণ Program লেখা
&
Script Mode
\\

\hline

Program Save ও Edit করা
&
Editor / Script Mode
\\

\hline

Project Development
&
IDE
\\

\hline

Data Analysis বা Experiment
&
Jupyter Notebook
\\

\hline

Installation ছাড়া Practice
&
Online Environment
\\

\hline

\end{tabularx}

%================================================
% First Script Program
%================================================

\booksubheading{একটি Script Program-এর উদাহরণ}

Editor-এ লিখি—

\begin{examplebox}{Program}

\begin{codebox}

\begin{english}
\ttfamily
print("Welcome to Python")
\end{english}

\end{codebox}

File Save করি—

\begin{codebox}

\begin{english}
\ttfamily
welcome.py
\end{english}

\end{codebox}

\textbf{Output:}

\begin{codebox}

\begin{english}
\ttfamily
Welcome to Python
\end{english}

\end{codebox}

\end{examplebox}

\textbf{কাজের ধারা—}

\begin{center}

\textbf{Editor}
\(\longrightarrow\)
\textbf{Code লেখা}
\(\longrightarrow\)
\textbf{Save as \textenglish{\ttfamily .py}}
\(\longrightarrow\)
\textbf{Run}
\(\longrightarrow\)
\textbf{Output}

\end{center}

%================================================
% Which one is important
%================================================

\booksubheading{Interactive Mode ও Script Mode—কখন কোনটি?}

দুটি Mode-ই গুরুত্বপূর্ণ; তবে ব্যবহার ভিন্ন।

\begin{itemize}

\item
দ্রুত কোনো Statement বা Calculation পরীক্ষা করতে
\(\longrightarrow\)
\textbf{Interactive Mode}

\item
পূর্ণ Program লিখে Save ও পুনরায় ব্যবহার করতে
\(\longrightarrow\)
\textbf{Script Mode}

\end{itemize}

\begin{rememberbox}

\textbf{\color{deepOrange}সহজ Shortcut}

\textbf{Practice / Quick Test}
\(\longrightarrow\)
Interactive Mode

\medskip

\textbf{Complete Program}
\(\longrightarrow\)
Script Mode

\end{rememberbox}

%================================================
% Common Mistakes
%================================================

\booksubheading{সাধারণ ভুল}

\begin{compacttwocolpoints}

\item
\textenglish{\ttfamily >>>}
-কে Python Program File-এর অংশ হিসেবে লিখে দেওয়া।

\item
Script File
\textenglish{\ttfamily .py}
Extension-এ Save না করা।

\item
Interactive Mode ও Script Mode-কে একই বিষয় মনে করা।

\item
Editor-কে Interpreter মনে করা।

\item
Code Editor এবং IDE-এর কাজ সম্পূর্ণ একই মনে করা।

\item
Shell-এ লেখা Code সবসময় স্বয়ংক্রিয়ভাবে Program File হিসেবে
Save হয় মনে করা।

\item
Interactive Mode-কে বড় Program লেখার প্রধান মাধ্যম মনে করা।

\end{compacttwocolpoints}

%================================================
% Final Reminder
%================================================

\begin{rememberbox}

\textbf{\color{deepOrange}মনে রাখবে}

\begin{itemize}

\item
Python Code প্রধানত
\textbf{Interactive Mode}
ও
\textbf{Script Mode}
ব্যবহার করে লেখা যায়।

\item
Interactive Mode-এ
\textenglish{\ttfamily >>>}
Prompt দেখা যায়।

\item
Interactive Mode-এ Instruction-এর Result দ্রুত পাওয়া যায়।

\item
Script Mode-এ Code Program File-এ লেখা ও Save করা হয়।

\item
Python Source File-এর Extension হলো
\textbf{\textenglish{\ttfamily .py}}।

\item
পূর্ণ Program লেখার জন্য Script Mode বেশি উপযোগী।

\item
IDLE-এ Python Shell ও Editor—দুইভাবেই কাজ করা যায়।

\item
VS Code Code Editor হিসেবে Python Development-এ ব্যবহার করা যায়।

\item
PyCharm একটি পরিচিত Python IDE।

\item
Jupyter Notebook-এ Code Cell-ভিত্তিকভাবে Execute করা যায়।

\item
দ্রুত মনে রাখবে—

\begin{center}

\textbf{Shell}
\(\longrightarrow\)
\textbf{Write}
\(\longrightarrow\)
\textbf{Enter}
\(\longrightarrow\)
\textbf{Result}

\medskip

\textbf{Script}
\(\longrightarrow\)
\textbf{Write}
\(\longrightarrow\)
\textbf{Save}
\(\longrightarrow\)
\textbf{Run}
\(\longrightarrow\)
\textbf{Output}

\end{center}

\end{itemize}

\end{rememberbox}

\small

%================================================
% 5.30 Python Program Introduction
%================================================

